\documentclass[11pt,a4paper]{article}
\usepackage[utf8]{inputenc}
\usepackage[T1]{fontenc}
\usepackage[catalan]{babel}
\usepackage{lmodern}
\usepackage[sfdefault,lining]{FiraSans}
\usepackage[a4paper,top=2.3cm,bottom=2.2cm,left=1.8cm,right=1.8cm,headheight=1.3cm,headsep=0.4cm,footskip=1cm]{geometry}
\usepackage{amsmath,amssymb}
\allowdisplaybreaks
\usepackage{icomma}
\usepackage{graphicx}
\usepackage[table]{xcolor}
\usepackage{tikz}
\usepackage{fontawesome5}
\usepackage{ccicons}
\usepackage{fancyhdr}
\usepackage{lastpage}
\usepackage{tasks}
\usepackage{needspace}
\usepackage[most]{tcolorbox}
\usepackage{array}
\usepackage[hidelinks,pdftitle={Solucions de la fitxa de la unitat 2: Sistemes d'equacions lineals. 2n de Batxillerat CCSS},pdfsubject={Matemàtiques Aplicades a les Ciències Socials II},pdfkeywords={sistemes d'equacions lineals, mètode de Gauss, mètode matricial, regla de Cramer, teorema de Rouché-Frobenius, sistemes homogenis, discussió de sistemes, paràmetres, problemes, PAU}]{hyperref}

\definecolor{navy}{HTML}{1B2748}
\definecolor{accent}{HTML}{E8603A}
\definecolor{teal}{HTML}{1F9E90}
\definecolor{soft}{HTML}{EEF1F8}
\definecolor{muted}{HTML}{5D6685}

\setlength{\parindent}{0pt}
\setlength{\parskip}{3pt}
\settasks{label=\alph*), label-format=\bfseries\color{navy}, label-width=1.3em, item-indent=1.8em, column-sep=0.6em, after-item-skip=5pt, before-skip=3pt, after-skip=4pt}

% ---------- capçalera i peu ----------
\newcommand{\llicencia}{\href{https://creativecommons.org/licenses/by-nc-sa/4.0/deed.ca}{\ccbyncsa}\ \ Llicència Creative Commons Reconeixement-NoComercial-CompartirIgual 4.0 Internacional (CC BY-NC-SA 4.0)}
\pagestyle{fancy}
\fancyhf{}
\fancyhead[L]{\raisebox{-0.35\height}{\includegraphics[height=0.95cm]{logo-gva.png}}}
\fancyhead[C]{\small\color{muted}Solucions de la fitxa · Unitat 2: Sistemes d'equacions lineals · 2n Batxillerat CCSS}
\fancyhead[R]{\raisebox{-0.35\height}{\includegraphics[height=0.95cm]{logo-vermellar.png}}}
\fancyfoot[L]{\scriptsize\color{muted}\llicencia}
\fancyfoot[R]{\small\color{muted}\thepage}
\renewcommand{\headrulewidth}{0.4pt}
\renewcommand{\headrule}{\color{muted!40}\hrule height \headrulewidth}
\fancypagestyle{primera}{\fancyhf{}\renewcommand{\headrulewidth}{0pt}
  \fancyfoot[L]{\scriptsize\color{muted}\llicencia}
  \fancyfoot[R]{\small\color{muted}\thepage}}

% ---------- nivells i exercicis ----------
\newcommand{\estrelles}[1]{{\color{accent}\ifcase#1\or\faStar\faStar[regular]\faStar[regular]\or\faStar\faStar\faStar[regular]\or\faStar\faStar\faStar\fi}}
\newcounter{exer}
\newcommand{\ex}[1]{\par\needspace{4\baselineskip}\medskip\refstepcounter{exer}\noindent{\color{navy}\bfseries\theexer.}\ \ }

\newcommand{\seccio}[2]{\par\needspace{10\baselineskip}\bigskip
  \begin{tcolorbox}[colback=navy,colframe=navy,coltext=white,arc=2mm,boxrule=0pt,left=3mm,right=3mm,top=1.2mm,bottom=1.2mm]
  {\large\bfseries #1\quad #2}
  \end{tcolorbox}\vspace{-2pt}}


% quadre de recordatori (només a la fitxa d'exercicis)
\newcommand{\recorda}[1]{\par\smallskip\begin{tcolorbox}[colback=accent!7,colframe=accent!55,arc=2mm,boxrule=0.5pt,left=2.5mm,right=2.5mm,top=1mm,bottom=1mm]\small{\color{accent}\faLightbulb[regular]}\ \textbf{Recorda:}\ #1\end{tcolorbox}}
% etiqueta dels exercicis de la PAU d'altres comunitats
\newcommand{\pau}[1]{{\color{teal}\small\bfseries\faUniversity\ PAU · #1.}\ }
\newcommand{\res}[1]{\boxed{#1}}\newcommand{\rest}[1]{\fbox{#1}}

% =====================================================================
% Capçalera FentMates de les fitxes: un color per a cada curs.
% Els colors són els de la pàgina inicial (LEVELS a assets/js/materials.js);
% el logo ve de _eines/logo/pdf/ (python3 _eines/logo/genera.py).
%
% Ús, al final del preàmbul de la unitat:
%   % =====================================================================
% Capçalera FentMates de les fitxes: un color per a cada curs.
% Els colors són els de la pàgina inicial (LEVELS a assets/js/materials.js);
% el logo ve de _eines/logo/pdf/ (python3 _eines/logo/genera.py).
%
% Ús, al final del preàmbul de la unitat:
%   % =====================================================================
% Capçalera FentMates de les fitxes: un color per a cada curs.
% Els colors són els de la pàgina inicial (LEVELS a assets/js/materials.js);
% el logo ve de _eines/logo/pdf/ (python3 _eines/logo/genera.py).
%
% Ús, al final del preàmbul de la unitat:
%   \input{../comu/fentmates.tex}
%   \fmcurs{1eso}
% i, al començament del document:
%   \capcalera{Matemàtiques}{1r ESO}{Fitxa d'exercicis · Unitat 2}{Divisibilitat}
% Opcional abans de \input: \def\fmneed{14} (espai mínim abans de cada secció).
% =====================================================================
\makeatletter
\def\fm@logodir{../../logo/pdf}
\providecommand{\fmneed}{10}
\def\fm@cursnom{}

% clau del curs, color i color del text damunt del color
\newcommand{\fm@curs}[3]{\@namedef{fm@color@#1}{#2}\@namedef{fm@text@#1}{#3}}
\fm@curs{1eso}{E6B84A}{navy}
\fm@curs{2eso}{3B82F6}{white}
\fm@curs{3eso}{FF7A50}{white}
\fm@curs{4eso-a}{A78BFA}{white}
\fm@curs{4eso-b}{6D28D9}{white}
\fm@curs{1batx-mi}{9B1C1C}{white}
\fm@curs{1batx-ccss}{EF6B6B}{white}
\fm@curs{2batx-mii}{1E6B47}{white}
\fm@curs{2batx-ccss}{2BB3A3}{white}

\newcommand{\fmcurs}[1]{%
  \@ifundefined{fm@color@#1}{\PackageError{fentmates}{Curs desconegut: #1}{}}{}%
  \edef\fm@tmp{\noexpand\definecolor{curs}{HTML}{\@nameuse{fm@color@#1}}%
               \noexpand\colorlet{curstext}{\@nameuse{fm@text@#1}}}\fm@tmp
  \colorlet{cursfosc}{curs!72!black}%
  \def\fm@logo{\fm@logodir/vertical-#1.pdf}%
  \def\fm@marca{\fm@logodir/marca-#1.pdf}}

% ---------- capçalera de la primera pàgina ----------
% #1 matèria (petit), #2 curs, #3 tipus de document i unitat, #4 títol
\newcommand{\capcalera}[4]{%
  \gdef\fm@cursnom{#2}%
  \par\noindent
  \begin{tcolorbox}[colback=curs,colframe=curs,boxrule=0pt,arc=2.2mm,width=3.5cm,height=2.15cm,
      valign=center,left=3mm,right=2mm,top=1mm,bottom=1mm,nobeforeafter,box align=center]
    \color{curstext}{\footnotesize #1}\par\vspace{1pt}{\fontsize{19}{21}\selectfont\bfseries #2\par}%
  \end{tcolorbox}\hspace{5mm}%
  \begin{minipage}[c]{\dimexpr\linewidth-3.5cm-5mm-2.9cm\relax}
    {\color{cursfosc}\small\bfseries\MakeUppercase{#3}}\par\vspace{2pt}
    {\color{navy}\fontsize{20}{23}\selectfont\bfseries #4\par}
  \end{minipage}\hfill
  \begin{minipage}[c]{2.6cm}\raggedleft\includegraphics[width=2.6cm]{\fm@logo}\end{minipage}\par
  \vspace{2pt}}

% ---------- barres de secció del color del curs ----------
\renewcommand{\seccio}[2]{\par\needspace{\fmneed\baselineskip}\bigskip
  \begin{tcolorbox}[colback=curs,colframe=curs,coltext=curstext,arc=2mm,boxrule=0pt,left=3mm,right=3mm,top=1.2mm,bottom=1.2mm]
  {\large\bfseries #1\quad #2}
  \end{tcolorbox}\vspace{-2pt}}

% ---------- capçalera i peu de les pàgines ----------
% primera pàgina: logos institucionals al peu
\fancypagestyle{primera}{\fancyhf{}\renewcommand{\headrulewidth}{0pt}
  \fancyfoot[L]{\raisebox{-0.32\height}{\includegraphics[height=0.78cm]{logo-gva.png}}\hspace{5mm}%
    \raisebox{-0.32\height}{\includegraphics[height=0.78cm]{logo-vermellar.png}}}
  \fancyfoot[R]{\scriptsize\color{muted}\href{https://creativecommons.org/licenses/by-nc-sa/4.0/deed.ca}{\ccbyncsa}\ \ CC BY-NC-SA 4.0 · fentmates.org\qquad\small\thepage}}
% la resta: el símbol a l'esquerra i el curs a la dreta, amb el títol al centre com abans
\fancyhead[L]{\raisebox{-0.3\height}{\includegraphics[height=0.62cm]{\fm@marca}}}
\fancyhead[R]{\small\bfseries\color{cursfosc}\fm@cursnom}
\renewcommand{\headrule}{\color{curs!45}\hrule height \headrulewidth}
\makeatother

%   \fmcurs{1eso}
% i, al començament del document:
%   \capcalera{Matemàtiques}{1r ESO}{Fitxa d'exercicis · Unitat 2}{Divisibilitat}
% Opcional abans de \input: \def\fmneed{14} (espai mínim abans de cada secció).
% =====================================================================
\makeatletter
\def\fm@logodir{../../logo/pdf}
\providecommand{\fmneed}{10}
\def\fm@cursnom{}

% clau del curs, color i color del text damunt del color
\newcommand{\fm@curs}[3]{\@namedef{fm@color@#1}{#2}\@namedef{fm@text@#1}{#3}}
\fm@curs{1eso}{E6B84A}{navy}
\fm@curs{2eso}{3B82F6}{white}
\fm@curs{3eso}{FF7A50}{white}
\fm@curs{4eso-a}{A78BFA}{white}
\fm@curs{4eso-b}{6D28D9}{white}
\fm@curs{1batx-mi}{9B1C1C}{white}
\fm@curs{1batx-ccss}{EF6B6B}{white}
\fm@curs{2batx-mii}{1E6B47}{white}
\fm@curs{2batx-ccss}{2BB3A3}{white}

\newcommand{\fmcurs}[1]{%
  \@ifundefined{fm@color@#1}{\PackageError{fentmates}{Curs desconegut: #1}{}}{}%
  \edef\fm@tmp{\noexpand\definecolor{curs}{HTML}{\@nameuse{fm@color@#1}}%
               \noexpand\colorlet{curstext}{\@nameuse{fm@text@#1}}}\fm@tmp
  \colorlet{cursfosc}{curs!72!black}%
  \def\fm@logo{\fm@logodir/vertical-#1.pdf}%
  \def\fm@marca{\fm@logodir/marca-#1.pdf}}

% ---------- capçalera de la primera pàgina ----------
% #1 matèria (petit), #2 curs, #3 tipus de document i unitat, #4 títol
\newcommand{\capcalera}[4]{%
  \gdef\fm@cursnom{#2}%
  \par\noindent
  \begin{tcolorbox}[colback=curs,colframe=curs,boxrule=0pt,arc=2.2mm,width=3.5cm,height=2.15cm,
      valign=center,left=3mm,right=2mm,top=1mm,bottom=1mm,nobeforeafter,box align=center]
    \color{curstext}{\footnotesize #1}\par\vspace{1pt}{\fontsize{19}{21}\selectfont\bfseries #2\par}%
  \end{tcolorbox}\hspace{5mm}%
  \begin{minipage}[c]{\dimexpr\linewidth-3.5cm-5mm-2.9cm\relax}
    {\color{cursfosc}\small\bfseries\MakeUppercase{#3}}\par\vspace{2pt}
    {\color{navy}\fontsize{20}{23}\selectfont\bfseries #4\par}
  \end{minipage}\hfill
  \begin{minipage}[c]{2.6cm}\raggedleft\includegraphics[width=2.6cm]{\fm@logo}\end{minipage}\par
  \vspace{2pt}}

% ---------- barres de secció del color del curs ----------
\renewcommand{\seccio}[2]{\par\needspace{\fmneed\baselineskip}\bigskip
  \begin{tcolorbox}[colback=curs,colframe=curs,coltext=curstext,arc=2mm,boxrule=0pt,left=3mm,right=3mm,top=1.2mm,bottom=1.2mm]
  {\large\bfseries #1\quad #2}
  \end{tcolorbox}\vspace{-2pt}}

% ---------- capçalera i peu de les pàgines ----------
% primera pàgina: logos institucionals al peu
\fancypagestyle{primera}{\fancyhf{}\renewcommand{\headrulewidth}{0pt}
  \fancyfoot[L]{\raisebox{-0.32\height}{\includegraphics[height=0.78cm]{logo-gva.png}}\hspace{5mm}%
    \raisebox{-0.32\height}{\includegraphics[height=0.78cm]{logo-vermellar.png}}}
  \fancyfoot[R]{\scriptsize\color{muted}\href{https://creativecommons.org/licenses/by-nc-sa/4.0/deed.ca}{\ccbyncsa}\ \ CC BY-NC-SA 4.0 · fentmates.org\qquad\small\thepage}}
% la resta: el símbol a l'esquerra i el curs a la dreta, amb el títol al centre com abans
\fancyhead[L]{\raisebox{-0.3\height}{\includegraphics[height=0.62cm]{\fm@marca}}}
\fancyhead[R]{\small\bfseries\color{cursfosc}\fm@cursnom}
\renewcommand{\headrule}{\color{curs!45}\hrule height \headrulewidth}
\makeatother

%   \fmcurs{1eso}
% i, al començament del document:
%   \capcalera{Matemàtiques}{1r ESO}{Fitxa d'exercicis · Unitat 2}{Divisibilitat}
% Opcional abans de \input: \def\fmneed{14} (espai mínim abans de cada secció).
% =====================================================================
\makeatletter
\def\fm@logodir{../../logo/pdf}
\providecommand{\fmneed}{10}
\def\fm@cursnom{}

% clau del curs, color i color del text damunt del color
\newcommand{\fm@curs}[3]{\@namedef{fm@color@#1}{#2}\@namedef{fm@text@#1}{#3}}
\fm@curs{1eso}{E6B84A}{navy}
\fm@curs{2eso}{3B82F6}{white}
\fm@curs{3eso}{FF7A50}{white}
\fm@curs{4eso-a}{A78BFA}{white}
\fm@curs{4eso-b}{6D28D9}{white}
\fm@curs{1batx-mi}{9B1C1C}{white}
\fm@curs{1batx-ccss}{EF6B6B}{white}
\fm@curs{2batx-mii}{1E6B47}{white}
\fm@curs{2batx-ccss}{2BB3A3}{white}

\newcommand{\fmcurs}[1]{%
  \@ifundefined{fm@color@#1}{\PackageError{fentmates}{Curs desconegut: #1}{}}{}%
  \edef\fm@tmp{\noexpand\definecolor{curs}{HTML}{\@nameuse{fm@color@#1}}%
               \noexpand\colorlet{curstext}{\@nameuse{fm@text@#1}}}\fm@tmp
  \colorlet{cursfosc}{curs!72!black}%
  \def\fm@logo{\fm@logodir/vertical-#1.pdf}%
  \def\fm@marca{\fm@logodir/marca-#1.pdf}}

% ---------- capçalera de la primera pàgina ----------
% #1 matèria (petit), #2 curs, #3 tipus de document i unitat, #4 títol
\newcommand{\capcalera}[4]{%
  \gdef\fm@cursnom{#2}%
  \par\noindent
  \begin{tcolorbox}[colback=curs,colframe=curs,boxrule=0pt,arc=2.2mm,width=3.5cm,height=2.15cm,
      valign=center,left=3mm,right=2mm,top=1mm,bottom=1mm,nobeforeafter,box align=center]
    \color{curstext}{\footnotesize #1}\par\vspace{1pt}{\fontsize{19}{21}\selectfont\bfseries #2\par}%
  \end{tcolorbox}\hspace{5mm}%
  \begin{minipage}[c]{\dimexpr\linewidth-3.5cm-5mm-2.9cm\relax}
    {\color{cursfosc}\small\bfseries\MakeUppercase{#3}}\par\vspace{2pt}
    {\color{navy}\fontsize{20}{23}\selectfont\bfseries #4\par}
  \end{minipage}\hfill
  \begin{minipage}[c]{2.6cm}\raggedleft\includegraphics[width=2.6cm]{\fm@logo}\end{minipage}\par
  \vspace{2pt}}

% ---------- barres de secció del color del curs ----------
\renewcommand{\seccio}[2]{\par\needspace{\fmneed\baselineskip}\bigskip
  \begin{tcolorbox}[colback=curs,colframe=curs,coltext=curstext,arc=2mm,boxrule=0pt,left=3mm,right=3mm,top=1.2mm,bottom=1.2mm]
  {\large\bfseries #1\quad #2}
  \end{tcolorbox}\vspace{-2pt}}

% ---------- capçalera i peu de les pàgines ----------
% primera pàgina: logos institucionals al peu
\fancypagestyle{primera}{\fancyhf{}\renewcommand{\headrulewidth}{0pt}
  \fancyfoot[L]{\raisebox{-0.32\height}{\includegraphics[height=0.78cm]{logo-gva.png}}\hspace{5mm}%
    \raisebox{-0.32\height}{\includegraphics[height=0.78cm]{logo-vermellar.png}}}
  \fancyfoot[R]{\scriptsize\color{muted}\href{https://creativecommons.org/licenses/by-nc-sa/4.0/deed.ca}{\ccbyncsa}\ \ CC BY-NC-SA 4.0 · fentmates.org\qquad\small\thepage}}
% la resta: el símbol a l'esquerra i el curs a la dreta, amb el títol al centre com abans
\fancyhead[L]{\raisebox{-0.3\height}{\includegraphics[height=0.62cm]{\fm@marca}}}
\fancyhead[R]{\small\bfseries\color{cursfosc}\fm@cursnom}
\renewcommand{\headrule}{\color{curs!45}\hrule height \headrulewidth}
\makeatother

\fmcurs{2batx-ccss}

\begin{document}
\thispagestyle{primera}

\capcalera{Matemàtiques CCSS}{2n Batx}{Solucions de la fitxa · Unitat 2}{Sistemes d'equacions lineals}
\vspace{8pt}

\begin{tcolorbox}[colback=soft,colframe=soft,arc=2mm,left=3mm,right=3mm,top=1.5mm,bottom=1.5mm]
\small Intenta resoldre cada exercici abans de mirar-ne la solució. Els passos de Gauss poden ser diferents dels teus i ser igualment correctes; el que ha de coincidir és la classificació i la solució (en un SCI, la solució general es pot escriure de maneres diferents segons la incògnita que prengues com a paràmetre). Pots comprovar-ho amb el resolutor de sistemes dels apunts interactius.
\end{tcolorbox}

% =====================================================
\seccio{1}{Sistemes d'equacions lineals. Mètode de Gauss}

\ex{1} \begin{tasks}(3)
  \task $5=5$, \ $-2=-2$, \ $7=7$: \rest{sí}
  \task La 2a: $2+1-2=1\neq-2$: \rest{no}
  \task La 1a: $0-1+3=2\neq5$: \rest{no}
\end{tasks}

\ex{1} \begin{tasks}(1)
  \task Sumant: $3x=3$, $x=1$, $y=2$. \rest{SCD}: rectes secants en $(1,2)$.
  \task La 2a és la 1a per $-2$: \rest{SCI}, la mateixa recta: $x=1+2\lambda,\ y=\lambda$.
  \task La 1a per 2 dona $6x+2y=4\neq5$: \rest{SI}, rectes paral·leles.
\end{tasks}

\ex{2} Hi ha moltes respostes possibles; per exemple:\begin{tasks}(1)
  \task La suma de les dues: \ $\res{3x+2z=5}$ \ (SCI).
  \task Mateixos coeficients que la 1a i un altre terme independent: \ $\res{x+y+z=0}$ \ (SI).
  \task \ $\res{z=1}$: \ queda $x+y=2,\ 2x-y=1$, \ d'on $x=1,\ y=1,\ z=1$ (SCD).
\end{tasks}

\ex{1} \begin{tasks}(1)
  \task \[A^*=\left(\begin{array}{rrr|r}1 & 1 & 1 & 6\\2 & -1 & 3 & 9\\1 & 2 & -1 & 2\end{array}\right)\xrightarrow[F_3'=F_3-F_1]{F_2'=F_2-2F_1}\left(\begin{array}{rrr|r}1 & 1 & 1 & 6\\0 & -3 & 1 & -3\\0 & 1 & -2 & -4\end{array}\right)\]\[\xrightarrow{F_2\leftrightarrow F_3}\left(\begin{array}{rrr|r}1 & 1 & 1 & 6\\0 & 1 & -2 & -4\\0 & -3 & 1 & -3\end{array}\right)\]\[\xrightarrow{F_3''=F_3'+3F_2'}\left(\begin{array}{rrr|r}1 & 1 & 1 & 6\\0 & 1 & -2 & -4\\0 & 0 & -5 & -15\end{array}\right)\]$-5z=-15\Rightarrow z=3$; \ $y-6=-4\Rightarrow y=2$; \ $x=6-2-3$: \ $\res{x=1;\ y=2;\ z=3}$
  \task \[A^*=\left(\begin{array}{rrr|r}2 & 1 & -1 & 1\\1 & -1 & 2 & 5\\3 & 2 & 1 & 10\end{array}\right)\xrightarrow{F_1\leftrightarrow F_2}\left(\begin{array}{rrr|r}1 & -1 & 2 & 5\\2 & 1 & -1 & 1\\3 & 2 & 1 & 10\end{array}\right)\]\[\xrightarrow[F_3'=F_3-3F_1]{F_2'=F_2-2F_1}\left(\begin{array}{rrr|r}1 & -1 & 2 & 5\\0 & 3 & -5 & -9\\0 & 5 & -5 & -5\end{array}\right)\]\[\xrightarrow{F_3''=3F_3'-5F_2'}\left(\begin{array}{rrr|r}1 & -1 & 2 & 5\\0 & 3 & -5 & -9\\0 & 0 & 10 & 30\end{array}\right)\]$10z=30\Rightarrow z=3$; \ $3y-15=-9\Rightarrow y=2$; \ $x=5+2-6$: \ $\res{x=1;\ y=2;\ z=3}$
\end{tasks}

\ex{2} \begin{tasks}(1)
  \task \[A^*=\left(\begin{array}{rrr|r}1 & -1 & 2 & 1\\2 & 1 & -1 & 3\\4 & -1 & 3 & 5\end{array}\right)\xrightarrow[F_3'=F_3-4F_1]{F_2'=F_2-2F_1}\left(\begin{array}{rrr|r}1 & -1 & 2 & 1\\0 & 3 & -5 & 1\\0 & 3 & -5 & 1\end{array}\right)\]\[\xrightarrow{F_3''=F_3'-F_2'}\left(\begin{array}{rrr|r}1 & -1 & 2 & 1\\0 & 3 & -5 & 1\\0 & 0 & 0 & 0\end{array}\right)\]\rest{SCI}. Fem $z=\lambda$: $3y=1+5\lambda$, \ $x=1+y-2\lambda$: \ $\res{x=\tfrac{4-\lambda}{3};\ y=\tfrac{1+5\lambda}{3};\ z=\lambda}$
  \task \[A^*=\left(\begin{array}{rrr|r}1 & -1 & 2 & 1\\2 & 1 & -1 & 3\\4 & -1 & 3 & 6\end{array}\right)\xrightarrow[F_3'=F_3-4F_1]{F_2'=F_2-2F_1}\left(\begin{array}{rrr|r}1 & -1 & 2 & 1\\0 & 3 & -5 & 1\\0 & 3 & -5 & 2\end{array}\right)\]\[\xrightarrow{F_3''=F_3'-F_2'}\left(\begin{array}{rrr|r}1 & -1 & 2 & 1\\0 & 3 & -5 & 1\\0 & 0 & 0 & 1\end{array}\right)\]L'última fila diu $0=1$: \ \rest{SI}, no té solució.
  \task La 2a i la 3a equacions són la 1a multiplicada per $-2$ i per $3$: només queda $x-2y+z=3$ (rang 1). \rest{SCI amb dos paràmetres}: \ $\res{x=3+2\lambda-\mu;\ y=\lambda;\ z=\mu}$
\end{tasks}

\ex{2} \pau{Navarra, EvAU extraordinària 2023}\[A^*=\left(\begin{array}{rrr|r}1 & 1 & 3 & 1\\4 & 3 & 5 & 5\\2 & 1 & -1 & 3\end{array}\right)\xrightarrow[F_3'=F_3-2F_1]{F_2'=F_2-4F_1}\left(\begin{array}{rrr|r}1 & 1 & 3 & 1\\0 & -1 & -7 & 1\\0 & -1 & -7 & 1\end{array}\right)\]\[\xrightarrow{F_3''=F_3'-F_2'}\left(\begin{array}{rrr|r}1 & 1 & 3 & 1\\0 & -1 & -7 & 1\\0 & 0 & 0 & 0\end{array}\right)\]Queda una fila de zeros: \rest{SCI}. Amb $z=\lambda$: \ $-y-7\lambda=1\Rightarrow y=-1-7\lambda$; \ $x=1-y-3\lambda=2+4\lambda$.\[\boxed{x=2+4\lambda;\ \ y=-1-7\lambda;\ \ z=\lambda,\quad\lambda\in\mathbb{R}}\]


% =====================================================
\seccio{2}{Expressió matricial. Mètode matricial}

\ex{2} \[\begin{pmatrix}2 & 0 & -1\\1 & 1 & 0\\0 & 1 & 1\end{pmatrix}\cdot\begin{pmatrix}x\\ y\\ z\end{pmatrix}=\begin{pmatrix}1\\ 0\\ 2\end{pmatrix}\]$|A|=\begin{vmatrix}2 & 0 & -1\\1 & 1 & 0\\0 & 1 & 1\end{vmatrix}=1\neq0$, per tant $A$ té inversa:\[\mathrm{Adj}(A)=\begin{pmatrix}1 & -1 & 1\\-1 & 2 & -2\\1 & -1 & 2\end{pmatrix},\qquad A^{-1}=\frac{1}{1}\big(\mathrm{Adj}(A)\big)^t=\begin{pmatrix}1 & -1 & 1\\-1 & 2 & -1\\1 & -2 & 2\end{pmatrix}\]\[X=A^{-1}\cdot B=\begin{pmatrix}1 & -1 & 1\\-1 & 2 & -1\\1 & -2 & 2\end{pmatrix}\begin{pmatrix}1\\ 0\\ 2\end{pmatrix}=\begin{pmatrix}3\\ -3\\ 5\end{pmatrix}\ \Rightarrow\ \boxed{x=3;\ y=-3;\ z=5}\]

\ex{1} $\res{\left\{\begin{aligned}x+2y&=4\\-x+3z&=1\end{aligned}\right.}$ \ \ Té \rest{2 equacions i 3 incògnites} (el producte d'una matriu $2\times3$ per una $3\times1$ és $2\times1$).\par No pot ser SCD: amb 2 equacions i 3 incògnites, en esglaonar sempre sobra almenys una incògnita. Les files de $A$ no són proporcionals i, amb $F_2'=F_2+F_1$, queda $2y+3z=5$: no apareix cap fila del tipus $0=d\neq0$. Per tant, és \rest{SCI}, amb un paràmetre: \ $x=-1+3\lambda;\ y=\tfrac{5-3\lambda}{2};\ z=\lambda$.

\ex{2} \begin{tasks}(1)
  \task $|A|=1$, \ $A^{-1}=\begin{pmatrix}1 & -2\\-1 & 3\end{pmatrix}$, \ $X=\begin{pmatrix}1 & -2\\-1 & 3\end{pmatrix}\begin{pmatrix}7\\ 3\end{pmatrix}=\begin{pmatrix}1\\ 2\end{pmatrix}$: \ $\res{x=1;\ y=2}$
  \task $|A|=2$, \ $A^{-1}=\frac12\begin{pmatrix}1 & 1 & -1\\-1 & 1 & 1\\1 & -1 & 1\end{pmatrix}$, \ $X=\frac12\begin{pmatrix}1 & 1 & -1\\-1 & 1 & 1\\1 & -1 & 1\end{pmatrix}\begin{pmatrix}3\\ 3\\ 4\end{pmatrix}=\frac12\begin{pmatrix}2\\ 4\\ 4\end{pmatrix}$: \ $\res{x=1;\ y=2;\ z=2}$
\end{tasks}

\ex{2} \pau{Galícia, ABAU extraordinària 2023}\[\begin{pmatrix}2 & 1 & 1\\0 & 2 & -1\\1 & 1 & 0\end{pmatrix}\cdot\begin{pmatrix}x\\ y\\ z\end{pmatrix}=\begin{pmatrix}2\\ -1\\ 1\end{pmatrix}\]$|A|=-1\neq0$, per tant $A$ té inversa:\[\mathrm{Adj}(A)=\begin{pmatrix}1 & -1 & -2\\1 & -1 & -1\\-3 & 2 & 4\end{pmatrix},\qquad A^{-1}=\frac{1}{-1}\big(\mathrm{Adj}(A)\big)^t=\begin{pmatrix}-1 & -1 & 3\\1 & 1 & -2\\2 & 1 & -4\end{pmatrix}\]\[X=A^{-1}\cdot B=\begin{pmatrix}-1 & -1 & 3\\1 & 1 & -2\\2 & 1 & -4\end{pmatrix}\begin{pmatrix}2\\ -1\\ 1\end{pmatrix}=\begin{pmatrix}2\\ -1\\ -1\end{pmatrix}\ \Rightarrow\ \boxed{x=2;\ y=-1;\ z=-1}\]

\ex{2} \pau{País Basc, PAU extraordinària 2025}\begin{tasks}(1)
  \task $\begin{pmatrix}3 & -2 & -2\\1 & 0 & -1\\0 & 2 & -1\end{pmatrix}\cdot\begin{pmatrix}x\\ y\\ z\end{pmatrix}=\begin{pmatrix}3\\ 1\\ 0\end{pmatrix}$
  \task $|A|=3\cdot2+(-2)\cdot1+(-2)\cdot2=6-2-4=0$: \ \rest{no té inversa} i no podem aplicar el mètode matricial.
  \task \[A^*=\left(\begin{array}{rrr|r}3 & -2 & -2 & 3\\1 & 0 & -1 & 1\\0 & 2 & -1 & 0\end{array}\right)\xrightarrow{F_1\leftrightarrow F_2}\left(\begin{array}{rrr|r}1 & 0 & -1 & 1\\3 & -2 & -2 & 3\\0 & 2 & -1 & 0\end{array}\right)\]\[\xrightarrow{F_2'=F_2-3F_1}\left(\begin{array}{rrr|r}1 & 0 & -1 & 1\\0 & -2 & 1 & 0\\0 & 2 & -1 & 0\end{array}\right)\]\[\xrightarrow{F_3'=F_3+F_2'}\left(\begin{array}{rrr|r}1 & 0 & -1 & 1\\0 & -2 & 1 & 0\\0 & 0 & 0 & 0\end{array}\right)\]\rest{SCI}. Amb $z=\lambda$: \ $2y=\lambda$, \ $x=1+z$: \ $\res{x=1+\lambda;\ y=\tfrac{\lambda}{2};\ z=\lambda,\ \ \lambda\in\mathbb{R}}$
\end{tasks}

\ex{3} \begin{tasks}(1)
  \task $|A|=\begin{vmatrix}1 & k & 0\\0 & 1 & 1\\k & 0 & 1\end{vmatrix}=1+k^2>0$ per a tot $k$: \ $A$ \rest{sempre és invertible}.
  \task $|A|=2$, \ $A^{-1}=\frac12\begin{pmatrix}1 & -1 & 1\\1 & 1 & -1\\-1 & 1 & 1\end{pmatrix}$, \ $X=\frac12\begin{pmatrix}1 & -1 & 1\\1 & 1 & -1\\-1 & 1 & 1\end{pmatrix}\begin{pmatrix}1\\ 2\\ 3\end{pmatrix}=\frac12\begin{pmatrix}2\\ 0\\ 4\end{pmatrix}$: \ $\res{x=1;\ y=0;\ z=2}$
\end{tasks}


% =====================================================
\seccio{3}{Regla de Cramer}

\ex{1} \begin{tasks}(1)
  \task $|A|=\begin{vmatrix}2 & -3\\5 & 1\end{vmatrix}=17$: \ $x=\dfrac{\begin{vmatrix}-4 & -3\\7 & 1\end{vmatrix}}{17}=\dfrac{17}{17}$, \ $y=\dfrac{\begin{vmatrix}2 & -4\\5 & 7\end{vmatrix}}{17}=\dfrac{34}{17}$: \ $\res{x=1;\ y=2}$
  \task $|A|=7$, \ $|A_x|=7$, \ $|A_y|=-7$, \ $|A_z|=14$: \ $\res{x=1;\ y=-1;\ z=2}$
\end{tasks}

\ex{2} \pau{Castella i Lleó, PAU ordinària 2025}Queda $\left\{\begin{aligned}x+2y+z&=1\\x-2y+2z&=4\\3x+y-2z&=5\end{aligned}\right.$, amb $|A|=\begin{vmatrix}1 & 2 & 1\\1 & -2 & 2\\3 & 1 & -2\end{vmatrix}=25\neq0$: és de Cramer.\[x=\frac{\begin{vmatrix}1 & 2 & 1\\4 & -2 & 2\\5 & 1 & -2\end{vmatrix}}{25}=\frac{52}{25},\quad y=\frac{\begin{vmatrix}1 & 1 & 1\\1 & 4 & 2\\3 & 5 & -2\end{vmatrix}}{25}=\frac{-17}{25},\quad z=\frac{\begin{vmatrix}1 & 2 & 1\\1 & -2 & 4\\3 & 1 & 5\end{vmatrix}}{25}=\frac{7}{25}\]\[\boxed{x=\tfrac{52}{25};\ \ y=-\tfrac{17}{25};\ \ z=\tfrac{7}{25}}\]

\ex{2} $|A|=\begin{vmatrix}1 & -1 & 2\\3 & 1 & -1\\2 & 3 & 1\end{vmatrix}=23\neq0$. \ $z=\dfrac{\begin{vmatrix}1 & -1 & 0\\3 & 1 & 2\\2 & 3 & 4\end{vmatrix}}{23}=\res{\dfrac{6}{23}}$; \ $x=\dfrac{10}{23}$, \ $y=\dfrac{22}{23}$ (i es comprova en les tres equacions).

\ex{3} $\begin{vmatrix}1&2\\ 2&-1\end{vmatrix}=-5\neq0$: \ $\left\{\begin{aligned}x+2y&=3+z\\ 2x-y&=1-3z\end{aligned}\right.$\[x=\frac{\begin{vmatrix}3+z&2\\ 1-3z&-1\end{vmatrix}}{-5}=\frac{-5+5z}{-5}=1-z,\qquad y=\frac{\begin{vmatrix}1&3+z\\ 2&1-3z\end{vmatrix}}{-5}=\frac{-5-5z}{-5}=1+z\]\[\boxed{x=1-\lambda;\ \ y=1+\lambda;\ \ z=\lambda,\quad\lambda\in\mathbb{R}}\]


% =====================================================
\seccio{4}{Sistemes homogenis}

\ex{1} \begin{tasks}(1)
  \task $|A|=\begin{vmatrix}1 & 2 & -1\\2 & -1 & 1\\3 & 1 & 1\end{vmatrix}=-5\neq0$: \ \rest{SCD}, només la solució trivial \ $\res{x=y=z=0}$.
  \task $|A|=0$ i $\begin{vmatrix}1&-1\\ 2&1\end{vmatrix}=3\neq0$: \ $\operatorname{rg}(A)=2<3$, \ \rest{SCI}. Sumant les dues primeres, $3x=0$; i de la 1a, $y=z$: \ $\res{x=0;\ y=\lambda;\ z=\lambda}$
\end{tasks}

\ex{2} \begin{tasks}(1)
  \task $|A|=\begin{vmatrix}1 & 1 & -2\\3 & -1 & 0\\-1 & 3 & -4\end{vmatrix}=0\Rightarrow\operatorname{rg}(A)<3$: \ \rest{SCI}, té infinites solucions a més de la trivial.
  \task \[A^*=\left(\begin{array}{rrr|r}1 & 1 & -2 & 0\\3 & -1 & 0 & 0\\-1 & 3 & -4 & 0\end{array}\right)\xrightarrow[F_3'=F_3+F_1]{F_2'=F_2-3F_1}\left(\begin{array}{rrr|r}1 & 1 & -2 & 0\\0 & -4 & 6 & 0\\0 & 4 & -6 & 0\end{array}\right)\]\[\xrightarrow{F_3''=F_3'+F_2'}\left(\begin{array}{rrr|r}1 & 1 & -2 & 0\\0 & -4 & 6 & 0\\0 & 0 & 0 & 0\end{array}\right)\]Amb $z=2\lambda$: \ $-4y+12\lambda=0\Rightarrow y=3\lambda$; \ $x=-y+2z=\lambda$: \ $\res{x=\lambda;\ y=3\lambda;\ z=2\lambda,\ \ \lambda\in\mathbb{R}}$
  \task $\lambda+3\lambda+2\lambda=6\lambda=12\Rightarrow\lambda=2$: \ $\res{x=2;\ y=6;\ z=4}$
\end{tasks}

\ex{1} És homogeni (compatible) i $\operatorname{rg}(A)\leq2<3$ incògnites: \rest{SCI}. Passant $z$ al segon membre: $2x-y=-z$, $x+3y=2z$ \ $\Rightarrow x=-\frac z7,\ y=\frac{5z}{7}$. Amb $z=7\lambda$: \ $\res{x=-\lambda;\ y=5\lambda;\ z=7\lambda}$


% =====================================================
\seccio{5}{Teorema de Rouché-Frobenius}

\ex{1} \label{ex:rf}\begin{tasks}(1)
  \task $\begin{vmatrix}1&2\\ 3&-1\end{vmatrix}=-7\neq0\Rightarrow\operatorname{rg}(A)=2$; \ $|A^*|=\begin{vmatrix}1 & 2 & 5\\3 & -1 & 1\\2 & 1 & 4\end{vmatrix}=0\Rightarrow\operatorname{rg}(A^*)=2=n$: \ \rest{SCD} (la solució és $x=1,\ y=2$).
  \task $\operatorname{rg}(A)=2$; \ $|A^*|=\begin{vmatrix}1 & 2 & 5\\3 & -1 & 1\\1 & 1 & 2\end{vmatrix}=7\neq0\Rightarrow\operatorname{rg}(A^*)=3$: \ \rest{SI}.
  \task $|A|=0$ i $\begin{vmatrix}1&1\\ 2&-1\end{vmatrix}=-3\neq0\Rightarrow\operatorname{rg}(A)=2$. Els menors d'ordre 3 de $A^*$ amb la columna de termes independents també valen 0 $\Rightarrow\operatorname{rg}(A^*)=2<3$: \ \rest{SCI} (la 3a equació és la 1a més el doble de la 2a).
\end{tasks}

\ex{2} \pau{Múrcia, PAU ordinària 2025}$|A|=\begin{vmatrix}3 & 2 & -1\\1 & -1 & 2\\2 & 3 & -1\end{vmatrix}=-10\neq0\Rightarrow\operatorname{rg}(A)=\operatorname{rg}(A^*)=3=n$: \ \rest{SCD}. Per Cramer o per Gauss: \ $\res{x=1;\ y=1;\ z=2}$

\ex{2} $\begin{vmatrix}1 & 1 & 1\\1 & -1 & 1\\2 & 1 & -1\end{vmatrix}=6\neq0\Rightarrow\operatorname{rg}(A)=3$; \ $|A^*|=0\Rightarrow\operatorname{rg}(A^*)=3=n$: \ \rest{SCD}. Suprimim la 4a equació (no forma part del menor) i resolem les tres primeres: \ $\res{x=1;\ y=1;\ z=1}$

\ex{2} \begin{tasks}(1)
  \task \rest{Falsa}: si $\operatorname{rg}(A)=\operatorname{rg}(A^*)=2$ és SCD (com al primer apartat de l'exercici \ref{ex:rf}).
  \task \rest{Certa}: $\operatorname{rg}(A)=3$ i $\operatorname{rg}(A^*)$ no pot passar de 3 (té 3 files): $\operatorname{rg}(A)=\operatorname{rg}(A^*)=3=n$.
  \task \rest{Falsa}: els rangs són diferents i el sistema és incompatible (cap solució).
  \task \rest{Falsa}: $\operatorname{rg}(A)\leq2<3$, així que, si és compatible, és indeterminat.
\end{tasks}


% =====================================================
\seccio{6}{Discussió de sistemes amb paràmetres}

\ex{2} \pau{Astúries, PAU extraordinària 2025}\begin{tasks}(1)
  \task $|A|=\begin{vmatrix}1&-1\\ 10&10+m\end{vmatrix}=20+m$. \ Té solució única (SCD) si i només si \ $\res{m\neq-20}$. (Per a $m=-20$: $x-y=20$ i $x-y=15$, SI.)
  \task Cal $|A|=0$: $\res{m=-20}$; aleshores la 2a equació és $x-y=15$ i, perquè coincidisca amb la 1a, $\res{\lambda=-15}$. Una solució: $\res{x=15,\ y=0}$ (en general, $x=15+t,\ y=t$).
\end{tasks}

\ex{2} \pau{Astúries, EBAU extraordinària 2024}\begin{tasks}(1)
  \task $A\cdot B=\begin{pmatrix}1&2m+1\\ -1-m&0\end{pmatrix}$, \ $A\cdot B-C=\begin{pmatrix}1&1\\ m-1&1\end{pmatrix}$: \ \ $\res{\left\{\begin{aligned}x+y&=1\\ (m-1)x+y&=1\end{aligned}\right.}$
  \task $|A|=2-m$. \ Si $m\neq2$: \rest{SCD}, solució única. \ Si $m=2$: les dues equacions són $x+y=1$, \rest{SCI}. \ Té solució \rest{per a tot $m$}, però només és única si $m\neq2$. \ Per a $m=-1$: $x+y=1$, $-2x+y=1$ \ $\Rightarrow\res{x=0;\ y=1}$.
\end{tasks}

\ex{2} \pau{Madrid, EvAU extraordinària 2022}\begin{tasks}(1)
  \task $|A|=\begin{vmatrix}1&a&1\\ 1&0&-a\\ 1&1&1\end{vmatrix}=1-a^2=0\Rightarrow a=\pm1$.\par $\bullet$ $a\neq\pm1$: \rest{SCD}. \ $\bullet$ $a=1$: la 1a i la 3a equacions són iguals, $\operatorname{rg}(A)=\operatorname{rg}(A^*)=2$: \rest{SCI}. \ $\bullet$ $a=-1$: $\operatorname{rg}(A)=2$ i $\begin{vmatrix}1&-1&2\\ 1&0&0\\ 1&1&2\end{vmatrix}=4\neq0\Rightarrow\operatorname{rg}(A^*)=3$: \rest{SI}.
  \task $x+z=2$, $x=0$, $x+y+z=2$: \ $\res{x=0;\ y=0;\ z=2}$
\end{tasks}

\ex{2} \pau{Madrid, EvAU extraordinària 2024}\begin{tasks}(1)
  \task $|A|=7a-21=0\Rightarrow a=3$. \ $\bullet$ $a\neq3$: \rest{SCD}. \ $\bullet$ $a=3$: $\begin{vmatrix}2&1\\ 3&1\end{vmatrix}=-1\neq0$, $\operatorname{rg}(A)=2$, i els menors d'ordre 3 de $A^*$ valen 0 (la 3a equació és $E_1+2E_2$): $\operatorname{rg}(A^*)=2$, \rest{SCI}.
  \task Restant $E_2-E_1$: $x-2z=-2$. Amb $z=\lambda$: $x=-2+2\lambda$, \ $y=-3x-z=6-7\lambda$:\par $\res{x=-2+2\lambda;\ y=6-7\lambda;\ z=\lambda}$
\end{tasks}

\ex{2} \pau{Madrid, EvAU ordinària 2024}\begin{tasks}(1)
  \task $|A|=2a^2-2a=2a(a-1)=0\Rightarrow a=0,\ a=1$.\par $\bullet$ $a\neq0,1$: \rest{SCD}. \ $\bullet$ $a=0$: $\operatorname{rg}(A)=2$ i $\begin{vmatrix}2&1&0\\ 1&0&1\\ 1&1&2\end{vmatrix}=-3\neq0$: \rest{SI}. \ $\bullet$ $a=1$: la 2a i la 3a equacions coincideixen ($x+y+z=2$), $\operatorname{rg}(A)=\operatorname{rg}(A^*)=2$: \rest{SCI}.
  \task $2x+y+z=1$ i $x+y+z=2$. Restant: $x=-1$; amb $z=\lambda$: \ $\res{x=-1;\ y=3-\lambda;\ z=\lambda}$
\end{tasks}

\ex{2} \pau{Madrid, PAU extraordinària 2025}\begin{tasks}(1)
  \task $|A|=2-a=0\Rightarrow a=2$. \ $\bullet$ $a\neq2$: \rest{SCD}. \ $\bullet$ $a=2$: la 3a equació, $3x+2z=5$, és la suma de les dues primeres: $\operatorname{rg}(A)=\operatorname{rg}(A^*)=2$, \rest{SCI}.
  \task Amb $z=\lambda$: $x+y=2-\lambda$, $2x-y=3-\lambda$ $\Rightarrow$ \ $\res{x=\tfrac{5-2\lambda}{3};\ y=\tfrac{1-\lambda}{3};\ z=\lambda}$
\end{tasks}

\ex{2} \pau{Múrcia, PAU extraordinària 2025}\begin{tasks}(1)
  \task $|A|=2-a=0\Rightarrow a=2$. \ $\bullet$ $a\neq2$: \rest{SCD}. \ $\bullet$ $a=2$: $\operatorname{rg}(A)=2$ i $\begin{vmatrix}1&-1&1\\ 2&-1&2\\ 3&2&4\end{vmatrix}=1\neq0\Rightarrow\operatorname{rg}(A^*)=3$: \rest{SI}.
  \task $E_2-E_1$: $x=0$. Aleshores $-y+z=1$ i $2y-z=4$: \ $\res{x=0;\ y=5;\ z=6}$
\end{tasks}

\ex{2} \pau{Múrcia, EBAU ordinària 2023}$|A|=-a^2+3a-2=-(a-1)(a-2)=0\Rightarrow a=1,\ a=2$.\par $\bullet$ $a\neq1,2$: \rest{SCD}. \ $\bullet$ $a=1$: la 1a i la 3a equacions tenen els mateixos coeficients i termes independents diferents: \rest{SI}. \ $\bullet$ $a=2$: la 1a i la 2a són iguals, $\operatorname{rg}(A)=\operatorname{rg}(A^*)=2$: \rest{SCI}.\par Per a $a=0$: $y+z=1$, $2x+y+z=1$, $x+z=2$ \ $\Rightarrow\res{x=0;\ y=-1;\ z=2}$

\ex{2} $|A|=\begin{vmatrix}1&1&-1\\ 2&-1&k\\ 1&2&1\end{vmatrix}=-k-8=0\Rightarrow\res{k=-8}$. Si $k\neq-8$: \rest{SCD}, només la solució trivial. Per a $k=-8$ és \rest{SCI}: de la 1a, $z=x+y$; a la 3a, $2x+3y=0$.\[\boxed{x=3\lambda;\ \ y=-2\lambda;\ \ z=\lambda,\quad\lambda\in\mathbb{R}}\]

\ex{3} \pau{Madrid, EvAU extraordinària 2023}\begin{tasks}(1)
  \task $|A|=a^3-7a+6=(a-1)(a-2)(a+3)=0\Rightarrow a=1,\ a=2,\ a=-3$.\par $\bullet$ $a\neq1,2,-3$: \rest{SCD}. \ $\bullet$ $a=1$: la 1a i la 2a equacions són iguals, $\operatorname{rg}(A)=\operatorname{rg}(A^*)=2$: \rest{SCI}. \ $\bullet$ $a=2$: la 2a i la 3a tenen els mateixos coeficients i termes independents diferents: \rest{SI}. \ $\bullet$ $a=-3$: $\operatorname{rg}(A)=2$ i $\begin{vmatrix}-3&1&1\\ 1&-3&-3\\ 1&2&1\end{vmatrix}=-8\neq0$: \rest{SI}.
  \task $y+2z=1$, $x+2z=0$, $x+2y=1$ \ $\Rightarrow\res{x=-\tfrac13;\ y=\tfrac23;\ z=\tfrac16}$
\end{tasks}

\ex{3} \pau{Castella i Lleó, PAU ordinària 2025}En general, $|A|=ab^2-2a-5b+13$ (l'exercici de la secció de Cramer resolia el cas $a=1$, $b=-2$).\begin{tasks}(1)
  \task $|A|=18-a=0\Rightarrow a=18$. \ $a\neq18$: \rest{SCD}. \ $a=18$: $\begin{vmatrix}1&-1\\ 3&1\end{vmatrix}=4\neq0$, $\operatorname{rg}(A)=2$, i $\begin{vmatrix}2&1&1\\ -1&2&4\\ 1&-1&5\end{vmatrix}=36\neq0$: \rest{SI}.
  \task $|A|=3b^2-5b+7$, que no s'anul·la mai (discriminant $25-84<0$): \rest{SCD per a tot $b$}.
\end{tasks}


% =====================================================
\seccio{7}{Problemes de plantejament}

\ex{1} \pau{Canàries, EBAU ordinària 2022}$x$, $y$, $z$: persones de pàdel surf, caiac i moto. \ \[\left\{\begin{aligned}x+y+z&=45\\40x+20y+60z&=1700\\x-3y&=0\end{aligned}\right.\] \ Amb $x=3y$: $4y+z=45$ i $7y+3z=85$\par $\Rightarrow$ \ \rest{30 de pàdel surf, 10 de caiac i 5 de moto aquàtica}.

\ex{2} \pau{Canàries, EBAU ordinària 2023}$x$: xiquets i joves, $y$: adults, $z$: jubilats. \ $\frac y4=\frac{2x}{5}$, \ $\frac{z}{x+y}=\frac{9}{26}$: \ \[\left\{\begin{aligned}x+y+z&=875\\-8x+5y&=0\\9x+9y-26z&=0\end{aligned}\right.\] \ De la 3a i la 1a: $35z=9\cdot875$, $z=225$; $x+y=650$ i $y=\frac85x$\par $\Rightarrow$ \ \rest{250 xiquets i joves, 400 adults i 225 jubilats}.

\ex{1} \pau{Canàries, EBAU extraordinària 2024}\[\left\{\begin{aligned}300x+400y+800z&=28000\\x-2y&=0\\y-2z&=0\end{aligned}\right.\] ($x$ mòbils, $y$ tauletes, $z$ portàtils). \ Amb $y=2z$, $x=4z$: $2800z=28\,000$\par $\Rightarrow$ \ \rest{40 mòbils, 20 tauletes i 10 portàtils}.

\ex{2} \pau{Canàries, EBAU ordinària 2024}$x$, $y$, $z$: tornejos locals, nacionals i internacionals. \ \[\left\{\begin{aligned}x-2y&=0\\2y-5z&=0\\0{,}3x+0{,}3y+0{,}3z&=51\end{aligned}\right.\] \ El total és $51/0{,}3=170$; amb $x=2y$ i $z=\frac25y$: $\frac{17}{5}y=170$\par $\Rightarrow$ \ \rest{100 locals, 50 nacionals i 20 internacionals}.

\ex{2} \pau{Cantàbria, EBAU extraordinària 2022}$x$, $y$, $z$: preu d'un gelat, d'un granissat i d'una orxata. Li han cobrat 4 gelats, 2 granissats i 3 orxates per 20 €; el que havia de pagar és 15 €; i el 50\,\% de $x+y$ és 2 €:\ \[\left\{\begin{aligned}4x+2y+3z&=20\\3x+2y+2z&=15\\x+y&=4\end{aligned}\right.\] \ Restant les dues primeres: $x+z=5$; amb $y=4-x$, $z=5-x$ en la 2a: $x=3$\par $\Rightarrow$ \ \rest{gelat 3 €, granissat 1 €, orxata 2 €}.

\ex{1} \pau{Cantàbria, EBAU ordinària 2023}Un sac de 25 kg val 100 €, un de 50 kg, 200 € i un de 100 kg, 400 €: \ \[\left\{\begin{aligned}x+y+z&=180\\100x+200y+400z&=29200\\x-2y-2z&=0\end{aligned}\right.\] \ De la 1a i la 3a: $3(y+z)=180$, $y+z=60$, $x=120$; i $y+2z=86$\par $\Rightarrow$ \ \rest{120 sacs de 25 kg, 34 de 50 kg i 26 de 100 kg}.

\ex{2} \pau{Cantàbria, EBAU extraordinària 2024}\begin{tasks}(1)
  \task $x$ banderes, $y$ samarretes, $z$ gorres: \ \[\left\{\begin{aligned}5x+6y+2z&=500\\2y-z&=0\\x+y&=70\end{aligned}\right.\]
  \task $|A|=\begin{vmatrix}5 & 6 & 2\\0 & 2 & -1\\1 & 1 & 0\end{vmatrix}=-5\neq0$: $\operatorname{rg}(A)=\operatorname{rg}(A^*)=3$, \rest{SCD}.
  \task \rest{40 banderes, 30 samarretes i 60 gorres}.
\end{tasks}

\ex{2} \pau{Cantàbria, EBAU ordinària 2024}\begin{tasks}(1)
  \task $x$, $y$, $z$: €/kg de terra, pedres i llavors en A. El proveïdor B cobra $13\,000-8\,800=4\,200$ €: \ \[\left\{\begin{aligned}300x+200y+100z&=13000\\100x+100y+20z&=4200\\2x+2y-z&=0\end{aligned}\right.\]
  \task $|A|=\begin{vmatrix}300 & 200 & 100\\100 & 100 & 20\\2 & 2 & -1\end{vmatrix}=-14000\neq0$: \rest{SCD}.
  \task Substituint $z=2x+2y$: $5x+4y=130$ i $x+y=30$\par $\Rightarrow$ \ \rest{terra 10 €/kg, pedres 20 €/kg, llavors 60 €/kg}.
\end{tasks}

\ex{1} \pau{Castella i Lleó, EBAU ordinària 2023}$x$ teatre, $y$ bàsquet, $z$ concert: \ \[\left\{\begin{aligned}0{,}9x+0{,}9y+0{,}9z&=117\\2x-z&=0\\-y+z&=20\end{aligned}\right.\] \ $x+y+z=130$; amb $x=\frac z2$, $y=z-20$: $\frac52z=150$\par $\Rightarrow$ \ \rest{teatre 30 €, bàsquet 40 €, concert 60 €}.

\ex{1} \pau{Castella-la Manxa, EvAU extraordinària 2022}$x$ Isabel, $y$ Carmen, $z$ Enma: $x+1=3z$, $z=\frac34y$: \ \[\left\{\begin{aligned}x+y+z&=15\\x-3z&=-1\\-3y+4z&=0\end{aligned}\right.\]\par $\Rightarrow$ \ \rest{Isabel 8, Carmen 4 i Enma 3}.

\ex{2} \pau{Castella-la Manxa, EvAU extraordinària 2022}\[\left\{\begin{aligned}10x+6y+3z&=851000\\x-y&=2000\\5x-6z&=13000\end{aligned}\right.\] \ Amb $y=x-2000$ i $z=\frac{5x-13000}{6}$ en la 1a: $18{,}5x=869\,500$.\par $\Rightarrow$ \ \rest{esportiu 47\,000 €, familiar 45\,000 €, monovolum 37\,000 €}.

\ex{1} \pau{Castella-la Manxa, EvAU ordinària 2022}\[\left\{\begin{aligned}12x+6y+6z&=51\\x-y&=1\\y-z&=-0{,}5\end{aligned}\right.\] \ Amb $x=y+1$, $z=y+0{,}5$: $24y+15=51$\par $\Rightarrow$ \ \rest{negra 2{,}50 €, amb llet 1{,}50 €, blanca 2 €}.

\ex{2} \pau{Castella-la Manxa, EvAU extraordinària 2023}\[\left\{\begin{aligned}x+y+z&=50\\2x-y&=1\\2x-5z&=0\end{aligned}\right.\] \ Amb $y=2x-1$, $z=\frac{2x}{5}$: $\frac{17}{5}x=51$\par $\Rightarrow$ \ \rest{15 de blanc, 29 de rosat i 6 de negre}.

\ex{2} \pau{Castella-la Manxa, EvAU ordinària 2023}$z-y=3(y-x)$: \ \[\left\{\begin{aligned}x+y+z&=70000\\x+y-z&=0\\3x-4y+z&=0\end{aligned}\right.\] \ Sumant les dues primeres, $z=35\,000$\par $\Rightarrow$ \ \rest{15\,000 del I, 20\,000 del II i 35\,000 del III}.

\ex{2} \pau{Castella-la Manxa, EvAU extraordinària 2024}$x$ Alba, $y$ Pablo, $z$ David: \ \[\left\{\begin{aligned}x+y+z&=72\\-x+y-z&=0\\x-\tfrac{1}{3}y-z&=0\end{aligned}\right.\] \ Sumant les dues primeres: $2y=72$, $y=36$\par $\Rightarrow$ \ \rest{24 Alba, 36 Pablo i 12 David}.

\ex{2} \pau{Castella-la Manxa, EvAU ordinària 2024}$x$ or, $y$ plata, $z$ bronze: $z+2=2(y-2)$: \ \[\left\{\begin{aligned}x+y+z&=36\\3x-z&=0\\2y-z&=6\end{aligned}\right.\]\par $\Rightarrow$ \ \rest{6 d'or, 12 de plata i 18 de bronze}.

\ex{3} \pau{Castella-la Manxa, EvAU ordinària 2024}$x$ Inés, $y$ Nerea, $z$ nuls: $4y-3x=z+1$, \ $x=y+7$, \ $z=0{,}05(x+y+z)\Leftrightarrow x+y-19z=0$: \ \[\left\{\begin{aligned}-3x+4y-z&=1\\x-y&=7\\x+y-19z&=0\end{aligned}\right.\]\par $\Rightarrow$ \ \rest{Inés 32 vots, Nerea 25 vots i 3 nuls}.

\ex{1} \pau{Castella-la Manxa, PAU extraordinària 2025}\[\left\{\begin{aligned}x+y+z&=25\\140x+130y+160z&=3610\\x-4y+z&=0\end{aligned}\right.\] \ De la 1a i la 3a: $5y=25$, $y=5$; després $7x+8z=148$ i $x+z=20$\par $\Rightarrow$ \ \rest{12 italians, 5 portuguesos i 8 japonesos}.

\ex{2} \pau{Catalunya, PAU extraordinària 2022}\begin{tasks}(1)
  \task $x$ autocar, $y$ a peu, $z$ taxi. El trajecte total és $x+y+z$: \ $x=\frac{x+y+z}{2}$, \ $y=\frac{x+y+z}{20}$, \ $x=z+5$:\ \[\left\{\begin{aligned}x-y-z&=0\\-x+19y-z&=0\\x-z&=5\end{aligned}\right.\] \ De la 1a i la 2a, $x=10y$ i $z=9y$; en la 3a, $y=5$: \ \rest{50 km en autocar, 5 km a peu i 45 km en taxi}.
  \task $\frac{50}{100}+\frac55+\frac{45}{90}=0{,}5+1+0{,}5=\res{2\text{ hores}}$.
\end{tasks}

\ex{3} \pau{Catalunya, PAU ordinària 2023}\begin{tasks}(1)
  \task $x$ termòmetre, $y$ mascareta, $z$ test: \ \[\left\{\begin{aligned}5x+30y+10z&=550\\x+15y+20z&=200\end{aligned}\right.\] \ Dues equacions i tres incògnites, $\operatorname{rg}(A)=\operatorname{rg}(A^*)=2<3$: \rest{SCI}, \textbf{no} es poden determinar els preus. Restant $5E_2-E_1$: $45y+90z=450$: \ $\res{y=10-2z;\ \ x=50+10z}$.
  \task $z=4$: \ \rest{la mascareta costa 2 € i el termòmetre, 90 €}.
\end{tasks}

\ex{1} \pau{Madrid, EvAU extraordinària 2024}$x$ guitarres, $y$ pianos, $z$ violins: \ \[\left\{\begin{aligned}x+y+z&=70\\x-y-z&=0\\4x+2y+z&=180\end{aligned}\right.\] \ De les dues primeres, $x=35$\par $\Rightarrow$ \ \rest{35 guitarres, 5 pianos i 30 violins}.

\ex{1} \pau{Madrid, EvAU ordinària 2024}\[\left\{\begin{aligned}x+y+z&=600\\10x+8y+5z&=4900\\x-2z&=0\end{aligned}\right.\] \ Amb $x=2z$: $3z+y=600$ i $25z+8y=4900$\par $\Rightarrow$ \ \rest{200 generals, 300 d'estudiant i 100 infantils}.

\ex{2} \pau{Navarra, PAU ordinària 2025}Les capses valen 6 €, 12 € i 24 €: \ \[\left\{\begin{aligned}x+y+z&=65\\x-y&=5\\6x+12y+24z&=720\end{aligned}\right.\]\par $\Rightarrow$ \ \rest{30 capses de 250 g, 25 de 500 g i 10 d'1 kg} (nombres naturals: la solució té sentit).

\ex{1} \pau{País Basc, PAU extraordinària 2025}\[\left\{\begin{aligned}x+y+z&=8\\10x+20y+50z&=290\\x-2y&=0\end{aligned}\right.\] \ Amb $x=2y$: $3y+z=8$ i $40y+50z=290$\par $\Rightarrow$ \ \rest{2 bitllets de 10 €, 1 de 20 € i 5 de 50 €}.

\ex{2} \pau{País Basc, PAU ordinària 2025}\[\left\{\begin{aligned}x+y+z&=20\\0{,}25x+0{,}5y+z&=10\\x+1{,}8y+3{,}3z&=35{,}6\end{aligned}\right.\] \ Multiplicant per 4 la 2a i per 10 la 3a, i restant la 1a: $y+3z=20$ i $8y+23z=156$\par $\Rightarrow$ \ \rest{8 llandes d'A, 8 de B i 4 de C}.

\ex{3} \pau{Astúries, PAU ordinària 2025}\begin{tasks}(1)
  \task $x+3y=90$, $2x+15y=225$\par $\Rightarrow$ \ \rest{VIP 75 €, normal 5 €}.
  \task Per a $m\neq6$ és SCD amb $x=\frac{45m}{m-6}$, $y=\frac{15(m-12)}{m-6}$. Si $6<m<12$, el preu $y$ és negatiu: per exemple, $\res{m=9}$ dona $x=135$, $y=-15$, que no té sentit.
  \task $|A|=m-6=0\Rightarrow\res{m=6}$: $2x+6y=90$ contradiu $x+3y=90$ (que multiplicada per 2 dona 180): \rest{SI}.
\end{tasks}

\ex{3} \pau{Astúries, EBAU extraordinària 2023}\begin{tasks}(1)
  \task $\res{\left\{\begin{aligned}x-my&=0\\ 0{,}15x+0{,}15y&=3m\end{aligned}\right.}$
  \task $|A|=0{,}15(1+m)$. Si $m\neq-1$: \rest{SCD} (solució única, $x=\frac{20m^2}{1+m}$, $y=\frac{20m}{1+m}$). Si $m=-1$: $x+y=0$ i $x+y=-20$: \rest{SI}. \ Amb 9 litres, $3m=9$, $m=3$: $x+y=60$, $x=3y$\par $\Rightarrow$ \ \rest{45 copes de vi negre i 15 de blanc}.
\end{tasks}


\end{document}
